\documentclass[11pt]{letter}
\usepackage[utf8]{inputenc}
\usepackage{hyperref}

% --- 用户需填写的信息 ---
\newcommand{\AuthorName}{Yuheng Fan}
\newcommand{\AuthorAffiliation}{Zhejiang Sci-Tech University}
\newcommand{\AuthorAddress}{Hangzhou, China}
\newcommand{\AuthorEmail}{2024110602002@mails.zstu.edu.cn}
\newcommand{\SubmissionDate}{\today}

% --- 期刊和编辑信息 ---
\newcommand{\EditorName}{Prof. Tulika Mitra}
\newcommand{\EditorTitle}{Editor-in-Chief}
\newcommand{\JournalName}{ACM Transactions on Embedded Computing Systems (TECS)}

% --- 稿件信息 ---
\newcommand{\ManuscriptTitle}{Efficient Verification Framework of LTL via Transition Certificates}

\begin{document}

% --- 发件人信息 ---
\begin{letter}{\EditorName \\ \EditorTitle \\ \JournalName}


\opening{Dear \EditorName,}

We are pleased to submit our original research article, titled "\textbf{\ManuscriptTitle}", for consideration for publication in \textbf{\JournalName}.

This manuscript presents a novel verification framework for Linear Temporal Logic (LTL) specifications on discrete-time dynamical systems, a critical challenge in the design of safe and reliable Cyber-Physical Systems (CPS). As the complexity of embedded systems grows, ensuring their correctness through formal verification becomes paramount. However, existing methods often face a trade-off between conservatism and high computational cost. Our work addresses this gap by introducing a more efficient approach.

The core innovation of our research is a new verification framework based on two novel types of certificates: \textit{transition safety certificates} and \textit{transition persistence certificates}. This approach effectively reduces the verification problem to analyzing the transitions of a Nondeterministic Büchi Automaton, significantly pruning the search space compared to traditional closure certificate methods while being less conservative than state-triplet based techniques. We also introduce an automated synthesis method based on a Counterexample-Guided Inductive Synthesis (CEGIS) framework, which supports both polynomial and neural network templates. Our experimental results on benchmark case studies, including the Kuramoto oscillator and a two-room temperature model, demonstrate a significant speedup over baseline methods, highlighting the practicality and scalability of our framework for real-world LTL verification tasks.

We believe this manuscript is an excellent fit for \textbf{\JournalName}. The journal's focus on the analysis, design, and behavior of embedded systems, particularly its interest in sophisticated algorithms, analytical models, and methodologies for CPS, aligns perfectly with the contributions of our paper. Our work on formal verification of LTL specifications directly addresses the challenges in ensuring the reliability of safety-critical embedded applications, a topic of great interest to your readership. We are confident that our findings will provide valuable insights and a practical solution for researchers and practitioners in the embedded systems community.

We confirm that this manuscript is original, has not been published elsewhere, and is not currently under consideration by any other journal. All authors have approved the manuscript and agree with its submission to \textbf{\JournalName}. The authors declare no competing interests.

Thank you for your time and consideration. We look forward to hearing from you.

\closing{Sincerely,\\
\vspace{1cm}
\AuthorName \\
\AuthorAffiliation \\
\AuthorAddress \\
\AuthorEmail \\
\SubmissionDate}

\end{letter}
\end{document}
